
%%%%%%%%%%%%%%%%%%%%%%%%%%%%%%%%%%%%%%%%%%%%%%%%
%%% Copier-Coller sur les compatibles IBM PC %%%
%%%%%%%%%%%%%%%%%%%%%%%%%%%%%%%%%%%%%%%%%%%%%%%%
\index{Copier}
\index{Coller}
\begin{figure}[htbp]
\begin{tabular}{@{}|p{13.2cm}|@{}} \hline
\multicolumn{1}{|c|}{\raisebox{-2ex}{\rule{0ex}{5ex}} \em Sélectionner et recopier un texte sur {\sc IBM PC} et compatibles }
\\  \hline \hline
{ % debut de la mini-page
\begin{minipage}[t]{13.2cm}
\vspace{2ex}
\begin{enumerate}
\item Placer le curseur d'édition au début du texte à 
recopier en le
déplaçant à l'aide des flèches de défilement
(\hspace{0.3ex}\begin{picture}(12,12)(0, 2)
\put(0,0){\framebox(12,12)[l]{$\:\scriptstyle\uparrow$}}
\end{picture}\hspace{0.3ex},
\begin{picture}(12,12)(0, 2)
\put(0,0){\framebox(12,12)[l]{$\:\scriptstyle\downarrow$}}
\end{picture}\hspace{0.3ex},
\begin{picture}(12,12)(0, 2)
\put(0,0){\framebox(12,12)[bl]{$\:\scriptstyle\leftarrow$}}
\end{picture}\hspace{0.3ex},
\begin{picture}(12,12)(0, 2)
\put(0,0){\framebox(12,12)[bl]{$\:\scriptstyle\rightarrow$}}
\end{picture}\hspace{0.3ex}).
\item Maintenir la touche de contrôle 
\begin{picture}(18,12)(0, 2)             %touche Controle
\put(0,0){\framebox(18,12)[lb]{}}
\put(1.5,2){\makebox(0,0)[lb]{\tiny Ctrl}}
\end{picture}\hspace{0.3ex} 
enfoncée le temps de taper 
\begin{picture}(12,12)(0, 2)             %touche K
\put(0,0){\framebox(12,12)[lb]{}}
\put(1.5,2){\makebox(0,0)[lb]{\tiny K}}
\end{picture}\hspace{0.3ex} 
suivi de \begin{picture}(12,12)(0, 2)    %touche B
\put(0,0){\framebox(12,12)[lb]{}}
\put(1.5,2){\makebox(0,0)[lb]{\tiny B}}
\end{picture}$\,$.
\item Placer le curseur à la fin du texte 
à recopier en le
déplaçant à l'aide des flèches de défilement
(\hspace{0.3ex}\begin{picture}(12,12)(0, 2)
\put(0,0){\framebox(12,12)[l]{$\:\scriptstyle\uparrow$}}
\end{picture}\hspace{0.3ex},
\begin{picture}(12,12)(0, 2)
\put(0,0){\framebox(12,12)[l]{$\:\scriptstyle\downarrow$}}
\end{picture}\hspace{0.3ex},
\begin{picture}(12,12)(0, 2)
\put(0,0){\framebox(12,12)[bl]{$\:\scriptstyle\leftarrow$}}
\end{picture}\hspace{0.3ex},
\begin{picture}(12,12)(0, 2)
\put(0,0){\framebox(12,12)[bl]{$\:\scriptstyle\rightarrow$}}
\end{picture}\hspace{0.3ex}), de sorte que le texte à recopier
apparaisse en {\em vidéo-inverse} c'est-à-dire en blanc sur fond noir.
\item Maintenir la touche de contrôle 
\begin{picture}(18,12)(0, 2)             %touche controle
\put(0,0){\framebox(18,12)[lb]{}}
\put(1.5,2){\makebox(0,0)[lb]{\tiny Ctrl}}
\end{picture}\hspace{0.3ex} 
enfoncée le temps de taper sur la touche
\begin{picture}(12,12)(0, 2)             %touche K
\put(0,0){\framebox(12,12)[lb]{}}
\put(1.5,2){\makebox(0,0)[lb]{\tiny K}}
\end{picture}\hspace{0.3ex} 
deux fois de suite.
\item Déplacer le curseur à l'aide des flèches de défilement
(\hspace{0.3ex}\begin{picture}(12,12)(0, 2)
\put(0,0){\framebox(12,12)[l]{$\:\scriptstyle\uparrow$}}
\end{picture}\hspace{0.3ex},
\begin{picture}(12,12)(0, 2)
\put(0,0){\framebox(12,12)[l]{$\:\scriptstyle\downarrow$}}
\end{picture}\hspace{0.3ex},
\begin{picture}(12,12)(0, 2)
\put(0,0){\framebox(12,12)[bl]{$\:\scriptstyle\leftarrow$}}
\end{picture}\hspace{0.3ex},
\mbox{\begin{picture}(12,12)(0, 2)
\put(0,0){\framebox(12,12)[bl]{$\:\scriptstyle\rightarrow$}}
\end{picture}\hspace{0.3ex})} à l'endroit où le texte 
doit être recopié.
\item Maintenir la touche de contrôle 
\begin{picture}(18,12)(0, 2)          %touche controle
\put(0,0){\framebox(18,12)[lb]{}}
\put(1.5,2){\makebox(0,0)[lb]{\tiny Ctrl}}
\end{picture}\hspace{0.3ex} 
enfoncée, puis taper 
\begin{picture}(12,12)(0, 2)          %touche K
\put(0,0){\framebox(12,12)[lb]{}}
\put(1.5,2){\makebox(0,0)[lb]{\tiny K}}
\end{picture}\hspace{0.3ex} 
suivi de \begin{picture}(12,12)(0, 2) %touche C
\put(0,0){\framebox(12,12)[lb]{}}
\put(1.5,2){\makebox(0,0)[lb]{\tiny C}}
\end{picture}$\,$.

\end{enumerate}
L'opération de recopie 
(\begin{picture}(18,12)(0, 2)          %touche controle
\put(0,0){\framebox(18,12)[lb]{}}
\put(1.5,2){\makebox(0,0)[lb]{\tiny Ctrl}}
\end{picture}\hspace{0.3ex} 
\begin{picture}(12,12)(0, 2)          %touche K
\put(0,0){\framebox(12,12)[lb]{}}
\put(1.5,2){\makebox(0,0)[lb]{\tiny K}}
\end{picture}\hspace{0.3ex} 
\begin{picture}(12,12)(0, 2)          %touche C
\put(0,0){\framebox(12,12)[lb]{}}
\put(1.5,2){\makebox(0,0)[lb]{\tiny C}}
\end{picture}\hspace{0.3ex}) en un endroit quelconque du
programme peut être répétée autant de
fois que nécessaire (sauf à l'intérieur de la zone recopiée).
\vspace{2ex}
\end{minipage} 
} % fin de la mini-page
\\ \hline
\end{tabular}
\caption{Marquage et copie de bloc sur les compatibles IBM PC}\label{fig:PC-copier-coller}
\end{figure}